\documentclass[class=scrbook, crop=false]{standalone}
\usepackage[subpreambles=true]{standalone}

\ifstandalone
    \input{../settings+/settings}
\fi

% ----------------------------------------------------------------------------
%                               Implementation
% ----------------------------------------------------------------------------
\begin{document}

\chapter{Primal Regularization for Indefinite Quadratic Programs}
\label{Chapter::Primal_Regularization}
In this chapter, we turn our attention to QPs with an indefinite Hessian matrix, namely indefinite QPs.
Section~\ref{sec:hessian_reg} classifies the existing primal regularization approaches by the matrix they target.
Section~\ref{sec:riccati_primal_reg} then develops the concept of Riccati-based primal regularization,
which applies the correction stagewise inside the Riccati recursion.
Section~\ref{sec:fail_minimal_perturbation} analyses why the minimum-perturbation rule fails in this setting,
and derives an improved stage regularization strategy from that analysis.
Section~\ref{sec:asymptotic_inactivity} concludes the chapter by proving that the proposed regularization
becomes asymptotically inactive as the iterates approach the solution.

\section{Primal regularization}
\label{sec:hessian_reg}
In indefinite QPs, the regularization needs to be applied to the Hessian matrix to ensure a descent search direction in the interior point method.
The resulting regularized reduced KKT system can be generally presented as:
\begin{equation}
    \begin{bmatrix}
     \tilde{H}_i + E_i & F^\top \\
    F & 0
    \end{bmatrix}
    \begin{bmatrix}
    \Delta y \\
    \Delta \pi
    \end{bmatrix} = -
    \begin{bmatrix}
    \tilde{r}_{\mathrm{stat}}^i \\
    r_{\mathrm{eq}}^i
    \end{bmatrix},
\end{equation}
where $E_i$ is the adjustment matrix that renders the Hessian positive definite at the $i$-th iteration.
This process is referred to as primal regularization.

Two choices are left open by this formulation: 
how the adjustment $E_i$ is computed, and which matrix it is derived from and applied to.
The two are independent, and the first is common ground,
since the same numerical strategies apply whichever matrix is targeted.
Section~\ref{subsec:reg_strategies} therefore introduces these strategies first,
after which Section~\ref{subsec:full_hessian} and Section~\ref{subsec:reduced_hessian}
distinguish the two types of primal regularization by their target matrix,
the full Hessian and the reduced Hessian.

\subsection{Regularization Strategies}
\label{subsec:reg_strategies}
Three families of strategies are commonly used to construct the adjustment matrix $E_i$
\cite{nocedal2006numerical, mcsweeneyModifiedCholeskyDecomposition2017, More1978levenberg}:
\begin{itemize}
    \item Eigenvalue modification (EM): replaces the negative eigenvalues of $\tilde{H}_i$ through a spectral decomposition.
    \item Modified Cholesky factorization (MC): corrects the inertia of $\tilde{H}_i$ by adjusting the diagonal pivots during factorization.
    \item Levenberg--Marquardt (LM): adds a scaled identity matrix to $\tilde{H}_i$ until it becomes positive definite.
\end{itemize}
The families differ in how the correction is obtained, and therefore in cost:
EM requires a full spectral decomposition, MC adjusts the pivots during the factorization itself,
and LM applies a uniform shift without any decomposition at all.
From these families we select four methods for this work:
PROJECT and MIRROR from EM \cite{verschueren2017sparsity},
GMW from MC \cite{gill1981practical},
and inertia correction from LM \cite{wachter2006implementation}.

Assume a matrix $A$ admits the eigenvalue decomposition $A = V D_E V^\top$
and the Cholesky factorization $A = L D_I L^\top$,
where $V$ is an orthogonal matrix of eigenvectors, $L$ is a unit lower triangular matrix,
$D_E$ is a diagonal matrix of real eigenvalues, and $D_I$ is a diagonal matrix whose signs give the inertia of $A$.
In this notation, the four methods read as follows:
\begin{itemize}
    \item \textbf{PROJECT} (Minimal eigenvalue modification):
    \begin{equation*}
        \mathrm{project}(A, \epsilon) := V \max(\epsilon I, D_E) V^\top.
    \end{equation*}
    Where the $\max$ operator is applied component-wise. 
    This approach clips any negative or near-zero eigenvalues to a small positive constant $\epsilon > 0$,
    which keeps a minimal adjustment to the original matrix for positive definiteness, thus the minimal Frobenius norm of the adjustment matrix $E$ \cite{nocedal2006numerical},
    but may lead to an increase in the condition number if the original eigenvalues were not close to zero.

    \item \textbf{MIRROR} (Absolute eigenvalue modification):
    \begin{equation*}
        \mathrm{mirror}(A, \epsilon) := V \max(\epsilon I, |D_E|) V^\top.
    \end{equation*}
    This method reflects negative eigenvalues to their positive counterparts. 
    It preserves the magnitude of the curvature in all directions, which often provides better scaling,
    but may result in a large modification to the original matrix.

    \item \textbf{GMW} (Modified Cholesky factorization):
    \begin{equation*}
        \mathrm{gmw}(A, \epsilon) := L D_I L^\top, \qquad D_I \succeq \epsilon I.
    \end{equation*}
    Instead of a spectral decomposition, the factorization is carried out in place, and each diagonal pivot is raised,
    whenever necessary, to keep it above $\epsilon$ and to bound the magnitude of the corresponding column of $L$ \cite{gill1981practical}.
    Bounding the pivots from below therefore also bounds the condition number of the factor,
    and the accumulated increments form a nonnegative diagonal correction $E$ with $A + E = L D_I L^\top$.

    This avoids the cost of an eigendecomposition and still guarantees positive definiteness, 
    but the resulting $E$ is generally not minimal and depends on the variable ordering, so it tends to over-modify $A$ compared with the eigenvalue-based methods.

    \item \textbf{Inertia Correction} (Levenberg--Marquardt):
    \begin{equation*}
        \mathrm{ic}(A, \Delta) := A + \Delta I, \qquad \Delta \geq 0.
    \end{equation*}
    A uniform multiple of the identity is added, so every eigenvalue is shifted by the same amount, $\lambda_k(A + \Delta I) = \lambda_k(A) + \Delta$.
    In practice $\lambda_{\min}$ is not computed directly, and $\Delta$ is found by repeatedly attempting the $\mathrm{LD}\mathrm{L}^\top$ factorization
    and increasing $\Delta$ until the inertia is correct, i.e.\ until every diagonal entry of $D_I$ is positive.
    The shift leaves the eigenvectors unchanged and is very cheap and robust,
    however, because it acts equally in all directions, it also inflates the already-positive eigenvalues and can degrade the scaling when only a few directions are indefinite.

\end{itemize}


\subsection{Full Hessian regularization}
\label{subsec:full_hessian}
The full hessian regularization take the full Hessian matrix $\tilde{H}$ as the target matrix.
Commonly, in the optimal control structure, the full Hessian is regularized stage-wisely 
by operating on the block Hessian matrix $\tilde{H}_k$ at $k$ stage,
\begin{equation}
    \label{eq:regularized_system}
    \tilde{H}_k^\mathrm{reg} =  
    \begin{bmatrix}
    \tilde{Q}_k & \tilde{S}_k^\top \\
    \tilde{S}_k & \tilde{R}_k
    \end{bmatrix} + E_k \succ 0, \quad k = 0, 1, \ldots, N-1,
\end{equation}
where $E_k$ could be a diagonal matrix or a full matrix, depending on the regularization strategy mentioned in \ref{subsec:reg_strategies}.

The full Hessian regularization is direct and usually simple to implement, 
but it can be too conservative without considering the definiteness of the reduced Hessian $\tilde{M}$, 
thus slow down the convergence.

\subsection{Reduced Hessian regularization}
\label{subsec:reduced_hessian}
The reduced Hessian regularization take the reduced Hessian matrix $\tilde{M}$ as the target matrix.
According to the Theorem \ref{thm:reduced_hessian_convexity}, 
which is mentioned in \cite{verschueren2017sparsity} (for proof, see \cite{nocedal2006numerical}):

\begin{theorem}
\label{thm:reduced_hessian_convexity}
Assume that LICQ holds. The QP problem has a unique global optimum if and only if the reduced Hessian is positive definite:
\begin{equation}
    M = Z^\top H Z \succ 0.
\end{equation}
\end{theorem}

The positive definiteness of the reduced Hessian is a necessary and sufficient condition for the QP admitting a unique global optimum.
It is thus reasonable to derive the correction from the reduced Hessian $\tilde{M}$ instead of full Hessian $\tilde{H}$, 
which yields a less conservative regularization that still effective.
However, from definition, 
computing the reduced Hessian directly requires a null-space basis of the equality constraints,
which can be achieved via condensing in optimal control problems \cite{frisonEfficientImplementationPartial2016}, 
this step generally introduces extra computational overhead.

We thus propose a Riccati-based primal regularization method, which naturally obtains the reduced Hessian $\tilde{M}$ implicitly via Riccati recursion,
and regularizes the reduced Hessian $\tilde{M}$ stage-wisely during the Riccati recursion to achieve a positive definite reduced Hessian $\tilde{M}$.

\section{Riccati-based Primal Regularization}
\label{sec:riccati_primal_reg}
In this section, we will introduce the idea and concept of Riccati-based primal regularization, 
a stage-wise regularization method for the reduced Hessian $\tilde{M}$ based on the Riccati recursion in two sense:
first, Riccati recursion is used for detecting the indefiniteness of the reduced Hessian $\tilde{M}$,
second, Riccati recursion is used for regularizing the reduced Hessian $\tilde{M}$
stage-wisely to ensure a positive definite reduced Hessian $\tilde{M}$.

\subsection{Block Reduced Hessian From Riccati Recursion}
\label{subsec:block_reduced_hessian}
As mentioned in \cite{frisonAlgorithmsMethodsHighPerformance}, 
the Riccati recursion can be interpreted as a block factorization of the KKT system in \eqref{eq:schur_comp_KKT}.
In this subsection, we will further show that the Riccati recursion 
can be regarded as the block $\mathrm{LD}\mathrm{L}^\top$ factorization of reduced Hessian $\tilde{M}_i$ in \eqref{eq:schur_comp_KKT_reduced},
which justifies the stage regularization on block reduced Hessian.

From the perspective of the Riccati recursion, the backward pass can be concluded as follows:
\begin{algorithm}[H]
\caption{Backward Pass of Riccati Recursion}
\label{alg:riccati_recursion_backward}
\begin{algorithmic}[1]
\State\textbf{Initialize:} $P_N \gets \tilde{Q}_N$, $p_N \gets \tilde{r}_{\mathrm{stat}, x}^N$
\Statex 
\For{$k = N-1$ \textbf{down to} $0$}
    \State $\tilde{M}_k \gets \tilde{R}_k + B_k^\top P_{k+1} B_k$
    \State $\Gamma_k \gets \tilde{S}_k + B_k^\top P_{k+1} A_k$
    \State $K_k \gets -\tilde{M}_k^{-1} \Gamma_k$
    \State $k_k \gets -\tilde{M}_k^{-1} \left(\tilde{r}_{\mathrm{stat}, u}^k + B_k^\top (P_{k+1} r_{\mathrm{eq}}^k + p_{k+1})\right)$
    \State $P_k \gets \tilde{Q}_k + A_k^\top P_{k+1} A_k + \Gamma_k^\top K_k$
    \State $p_k \gets \tilde{r}_{\mathrm{stat}, x}^k + A_k^\top (P_{k+1} r_{\mathrm{eq}}^k + p_{k+1}) + \Gamma_k^\top k_k$
\EndFor
\end{algorithmic}
\end{algorithm}
We can find that $\tilde{M}_k = \tilde{R}_k + B_k^\top P_{k+1} B_k$ is exactly the block reduced Hessian obtained by Riccati recursion.
In practice the inverse of $\tilde{M}_k$ is not computed directly; instead, the Cholesky factorization of $\tilde{M}_k$ is computed for efficiency.

From the perspective of the reduced Hessian $\tilde{M}_i$ in \eqref{eq:schur_comp_KKT_reduced}, 
eliminating the states through the dynamics leaves a dense reduced Hessian $\tilde{M}_i$ in the inputs $u_0, \dots, u_{N-1}$, 
whose density stems from the coupling of any two stages via the state transition $\Phi_{i, j}$, defined as:
\begin{align*}
    \Phi_{i, j} = \begin{cases}
    \prod_{k=j}^{i-1} A_k = A_{i-1} A_{i-2} \cdots A_j, & \text{if } i > j, \\
    I, & \text{if } i = j.
    \end{cases}
\end{align*}

The backward pass in Algorithm~\ref{alg:riccati_recursion} is exactly a block $\mathrm{LD}\mathrm{L}^\top$ factorization of this $\tilde{M}_i$, 
carried out without ever forming it explicitly. The connection is made through the Schur complement: 
each cost-to-go matrix $P_k$ is the Schur complement obtained by eliminating the states of the later stages, 
so that the pivot block $\tilde{M}_k$ is precisely the Schur complement of stage $k$ in $\tilde{M}_i$ after all later stages have been eliminated. 
Consequently, the recursion $P_k = \tilde{Q}_k + A_k^\top P_{k+1} A_k - \Gamma_k^\top \tilde{M}_k^{-1} \Gamma_k$ is nothing but the successive formation of these Schur complements.

Collecting the pivots $\tilde{M}_{N-1}, \dots, \tilde{M}_0$ into the block-diagonal factor $D$ and eliminating the stages in the same order as the backward pass, 
from stage $N-1$ down to stage $0$, the block $\mathrm{LD}\mathrm{L}^\top$ factorization of the reduced Hessian $\tilde{M}_i$ can be expressed as:
\begin{align*}
    \tilde{M}_i &= L
    \begin{bmatrix}
    \tilde{M}_{N-1} & 0 & \cdots & 0 \\
    0 & \tilde{M}_{N-2} & \cdots & 0 \\
    \vdots & \vdots & \ddots & \vdots \\
    0 & 0 & \cdots & \tilde{M}_{0}
    \end{bmatrix}
    L^\top
\end{align*}
where the accumulated back-substitution coefficients, built from the feedback gains $K_k$ and the transitions $\Phi_{i, j}$, form the unit lower-triangular factor $L$:
\begin{align*}
    L = 
    \begin{bmatrix}
    I & 0 & \cdots & 0 & 0 \\
    -B_{N-2}^\top K_{N-1}^\top & I & \cdots & 0 & 0 \\
    -B_{N-3}^\top \Phi_{N-1, N-2}^\top K_{N-1}^\top & -B_{N-3}^\top K_{N-2}^\top & \cdots & 0 & 0 \\
    \vdots & \vdots & \ddots & \vdots & \vdots \\
    -B_0^\top \Phi_{N-1, 1}^\top K_{N-1}^\top & -B_0^\top \Phi_{N-2, 1}^\top K_{N-2}^\top & \cdots & -B_0^\top K_1^\top & I
    \end{bmatrix}.
\end{align*}

This establishes the connection between positive definiteness of the reduced Hessian $\tilde{M}_i$ and the positive definiteness of the block reduced Hessians $\tilde{M}_k$ obtained from Riccati recursion, 
which justifies the stage regularization on block reduced Hessian $\tilde{M}_k$.

\subsection{Stage Regularization in Riccati Recursion}
With this connection, we can regularize the reduced Hessian $\tilde{M}_i$ without explicitly computing it, 
but by stage-wisely applying the regularization to the block reduced Hessian  $\tilde{M}_k^i$ during the Riccati recursion,

As mentioned in \ref{subsec:block_reduced_hessian}, 
the Cholesky factorization of $\tilde{M}_k^i$ is computed for efficiency, which requires the positive definiteness of $\tilde{M}_k^i$.
This can be regarded as a detection of the indefiniteness of the reduced Hessian $\tilde{M}_i$.
And when the Cholesky factorization fails, which indicates that the reduced Hessian $\tilde{M}_i$ is indefinite,
we can regularize the reduced Hessian block $\tilde{M}_k^i$ by adding a block adjust matrix $E_k^i$, i.e. $\tilde{M}_k^i + E_k^i$.
The $E_k^i$ can be computed from the strategies mentioned in \ref{subsec:reg_strategies}.
This gives a modification of the backward pass of Riccati recursion and can be concluded as follows:
\begin{algorithm}[H]
\caption{Backward Pass of Modified Riccati Recursion}
\label{alg:modified_riccati_recursion}
\begin{algorithmic}[1]
\State\textbf{Initialize:} $P_N \gets \tilde{Q}_N$, $p_N \gets \tilde{r}_{\mathrm{stat}, x}^N$, $\lambda > 0$
\Statex 
\For{$k = N-1$ \textbf{down to} $0$}
    \State $\tilde{M}_k \gets \tilde{R}_k + B_k^\top P_{k+1} B_k$
    \State $L_k \gets \mathrm{Cholesky}(\tilde{M}_k)$
    \If{factorization fails}
        \State $\tilde{M}_k \gets \tilde{M}_k + E_k$ \Comment{Apply regularization}
        \State $L_k \gets \mathrm{Cholesky}(\tilde{M}_k)$
    \EndIf
    \State $\Gamma_k \gets \tilde{S}_k + B_k^\top P_{k+1} A_k$
    \Statex
    \State $K_k \gets -L_k^{-\top} L_k^{-1} \Gamma_k$
    \State $k_k \gets -L_k^{-\top} L_k^{-1} \left(\tilde{r}_{\mathrm{stat}, u}^k + B_k^\top (P_{k+1} r_{\mathrm{eq}}^k + p_{k+1})\right)$
    \Statex
    \State $P_k \gets \tilde{Q}_k + A_k^\top P_{k+1} A_k + \Gamma_k^\top K_k$
    \State $p_k \gets \tilde{r}_{\mathrm{stat}, x}^k + A_k^\top (P_{k+1} r_{\mathrm{eq}}^k + p_{k+1}) + \Gamma_k^\top k_k$
\EndFor
\end{algorithmic}
\end{algorithm}
Since the block reduced Hessians $\tilde{M}_k$ are exactly the pivots of
the block $\mathrm{L}^\top \mathrm{DL}$ factorization of $\tilde{M}_i$,
enforcing the positive definiteness of every $\tilde{M}_k$ during the backward pass is equivalent to regularizing the reduced Hessian $\tilde{M}_i$ itself, 
and thus recovers the SOSC in Theorem~\ref{thm:reduced_hessian_convexity}.
The regularization is therefore applied only where it is needed while the structure-exploiting complexity of the Riccati recursion is preserved.


\section{Improved Stage Regularization Strategy}
\label{sec:fail_minimal_perturbation}
In this section, we will analyze the case 
where the common-used minimum perturbation strategy in regularization fails in stage regularization,
which causes a severe condition number issue in the preceding stages due to the propagation of the negative curvature through the stages.
This motivates us to propose an improved stage regularization strategy that avoids minimum perturbation and takes the propagation into account.

\subsection{Failure of Minimum-Perturbation Regularization}
In this subsection, we will show that the minimum perturbation strategy, i.e. the PROJECT \cite{verschueren2017sparsity}, 
which is commonly the first choice for regularization with the least modification of the original problem, 
can cause condition number issue in the stage regularization by introducing a general example.

Suppose at the stage $k$, the reduced Hessian block $\tilde{M}_k$ is not positive definite.
Without loss of generality, we can assume $\tilde{M}_k$ has only one negative eigenvalue $\lambda_j$, and the corresponding eigenvector is $v_j$.
The minimum perturbation strategy is applied as:
\begin{equation}
    E_k = (-\lambda_{j}+\epsilon) v_j v_j^\top,
\end{equation}
where the $\epsilon$ is a small positive value to ensure the positive definiteness of the regularized reduced Hessian block. 

By Sherman-Morrison formula, the inverse of the regularized reduced Hessian block can be expressed as:
\begin{align*}
    (\tilde{M}_k + E_k)^{-1} &= \tilde{M}_k^{-1} + \frac{\lambda_j-\epsilon}{\epsilon \lambda_j} v_j v_j^\top \\
                            &= \tilde{M}_k^{-1} + (\frac{1}{\epsilon} - \frac{1}{\lambda_j}) v_j v_j^\top,
\end{align*}
and the update of cost-to-go matrix $P_k$ follows as:
\begin{align}
    P_k &= \tilde{Q}_k + A_k^\top P_{k+1} A_k - \Gamma_k^\top K_k \notag \\
        &=\tilde{Q}_k + A_k^\top P_{k+1} A_k - \Gamma_k^\top \tilde{M}_k^{-1} \Gamma_k - (\frac{1}{\epsilon} - \frac{1}{\lambda_j}) \Gamma_k^\top v_j v_j^\top \Gamma_k \notag \\
        &= P_k^{\mathrm{unreg}} - (\frac{1}{\epsilon} - \frac{1}{\lambda_j}) \Gamma_k^\top v_j v_j^\top \Gamma_k,
\end{align}
note that this negative rank-1 modification can already render $P_k$ itself indefinite, 
independently of its downstream effect on $\tilde{M}_{k-1}$, 
whenever the perturbation magnitude exceeds the smallest eigenvalue of $P_k^{\mathrm{unreg}}$. 
This can happen even when the perturbation on $\tilde{M}_k$ is moderate,
provided that the negative-curvature direction $v_j$ is strongly coupled to the state through $\Gamma_k$.

Consequently, the reduced Hessian block $\tilde{M}_{k-1}$ at the previous stage as:
\begin{equation}
    \tilde{M}_{k-1} = \tilde{M}_{k-1}^{\mathrm{unreg}} - (\frac{1}{\epsilon} - \frac{1}{\lambda_j}) z_j z_j^\top,
\end{equation}
where $z_j = B_{k-1}^\top \Gamma_k^\top v_j$ is a column vector.

And the condition number of the regularized reduced Hessian block $\tilde{M}_{k-1}$ can be approximated as:
\begin{equation}
    \kappa(\tilde{M}_{k-1}) \approx \frac{c (\epsilon + |\lambda_j|)}{\epsilon |\lambda_j|},
\end{equation}
where $c = \frac{\|z_j\|^2}{\lambda_{\min}(\tilde{M}_{k-1}^{\mathrm{unreg}})}$ as an $O(1)$ constant.

Notably, as the $\epsilon$ is a small positive value,
the presence of the negative rank-1 update $-\frac{1}{\epsilon} z_j z_j^\top$ will dominate, not only severely deteriorates its condition number,
but also likely to make $\tilde{M}_{k-1}$ indefinite, which can further cause a domino effect through the stages.

\subsection{Insights and Proposed Methods}
This issue reveals two insights in stage regularization:
\begin{itemize}
    \item The minimum modification objective \cite{fangModifiedCholeskyAlgorithms2007} is not the foremost consideration for stage regularization, 
    as it can cause severe condition number issue in the preceding stages.
    A trade-off between perturbation norm and condition number must be carefully treated.
    \item The cost-to-go matrix $P_k$ must be incorporated into the regularization design, 
    as the severity of propagation depends on the coupling structure between stages, 
    which a purely local scheme cannot capture.
\end{itemize}

From the first insight, instead of the minimum perturbation strategy,
an intuitive, but effective, approach is to use the absolute value of the negative eigenvalues or inertia as the perturbation magnitude,
which can avoid the severe condition number issue in the preceding stages.
When it comes to the eigenvalue modification strategy, this is the MIRROR strategy \cite{verschueren2017sparsity},
and when it comes to the modified Cholesky factorization strategy, this is the Type-II strategy \cite{fangModifiedCholeskyAlgorithms2007}, of which GMW is a common representitive.

From the second insight, we can incorporate the cost-to-go matrix $P_k$ into the strategy design,
which also can be achieved by adding an adjustment matrix $E_{k+1,P}^i$ to the cost-to-go matrix $P_k$ directly,
so that the propagation of the negative curvature can be taken into account in the regularization design.

The resulting improved stage regularization strategy can be summarized as follows:

\begin{algorithm}[H]
\caption{Backward Pass of Improved Modified Riccati Recursion}
\label{alg:improved_modified_riccati_recursion}
\begin{algorithmic}[1]
\State\textbf{Initialize:} $P_N \gets \tilde{Q}_N$, $p_N \gets \tilde{r}_{\mathrm{stat}, x}^N$, $\mathcal{E}(\cdot)$: MIRROR or Type-II strategy
\Statex 
\For{$k = N-1$ \textbf{down to} $0$}
    \State $P_{k+1}^\mathrm{reg} \gets P_{k+1}$ \Comment{Default with no regularization}
    \State $\tilde{M}_k \gets \tilde{R}_k + B_k^\top P_{k+1}^\mathrm{reg} B_k$
    \State $L_k \gets \mathrm{Cholesky}(\tilde{M}_k)$
    \If{factorization fails}
        \State $P_{k+1}^\mathrm{reg} \gets \mathcal{E}(P_{k+1})$ \Comment{Regularize cost-to-go matrix}
        \State $\tilde{M}_k \gets \tilde{R}_k + B_k^\top P_{k+1}^\mathrm{reg} B_k$
        \State $\tilde{M}_k^\mathrm{reg} \gets \mathcal{E}(\tilde{M}_k)$ \Comment{Regularize reduced Hessian block}
        \State $L_k \gets \mathrm{Cholesky}(\tilde{M}_k^\mathrm{reg})$
    \EndIf
    \State $\Gamma_k \gets \tilde{S}_k + B_k^\top P_{k+1}^\mathrm{reg} A_k$
    \Statex
    \State $K_k \gets -L_k^{-\top} L_k^{-1} \Gamma_k$
    \State $k_k \gets -L_k^{-\top} L_k^{-1} \left(\tilde{r}_{\mathrm{stat}, u}^k + B_k^\top (P_{k+1}^\mathrm{reg} r_{\mathrm{eq}}^k + p_{k+1})\right)$
    \Statex
    \State $P_k \gets \tilde{Q}_k + A_k^\top P_{k+1}^\mathrm{reg} A_k + \Gamma_k^\top K_k$
    \State $p_k \gets \tilde{r}_{\mathrm{stat}, x}^k + A_k^\top (P_{k+1}^\mathrm{reg} r_{\mathrm{eq}}^k + p_{k+1}) + \Gamma_k^\top k_k$
\EndFor
\end{algorithmic}
\end{algorithm}

In Algorithm~\ref{alg:improved_modified_riccati_recursion}, when the Cholesky factorization fails, the regularization acts at two levels. 
Applying $\mathcal{E}(\cdot)$ to the cost-to-go matrix proactively convexifies the term that enters the current stage, 
so that the negative-curvature propagation identified in the second insight is accounted for in $\tilde{M}_k$, $\Gamma_k$, and $P_k$ consistently. 
Applying it again to the reduced Hessian block acts as a local safeguard, ensuring that the Cholesky factorization succeeds.
Both corrections are produced by the same non-minimal strategy $\mathcal{E}(\cdot)$, MIRROR or Type-II, rather than by the minimum-modification rule, 
which keeps the condition number of the regularized blocks bounded, as motivated by the first insight. 
When no factorization failure occurs, the recursion reduces to the standard Riccati backward pass.


\section{Asymptotic Inactivity of the Regularization}
\label{sec:asymptotic_inactivity}
In this section,
we will prove that the reduced Hessian regularization via Riccati recursion is asymptotically inactive,
which means that the regularization will be inactive when the iterates are sufficiently close to the optimal solution $(y^*, \lambda^*)$,
following a similar proof strategy to \cite{absilNewtonKKTInteriorpointMethods2007}.

We consider the regularized KKT formulation, and apply the null space method to \eqref{eq:IPM_KKT_system} first.
The resulting formulation is given as:
\begin{equation}
    \begin{bmatrix}
    W_i & Z^\top G^\top\\
    GZ & -\Lambda_i^{-1}T_i\\
    \end{bmatrix}
    \begin{bmatrix}
    \Delta y_z \\
    \Delta \lambda \\
    \end{bmatrix} = -
    \begin{bmatrix}
    r_{\mathrm{stat, redu}}^i \\
    \tilde{r}_{\mathrm{ineq, redu}}^i \\
    \end{bmatrix},
    \label{eq:reduced_KKT_system_schur}
\end{equation}
where
\begin{align*}
    W_i &= Z^\top H Z + E_i = M + E_i, \\
    r_{\mathrm{stat, redu}}^i &= Z^\top r_{\mathrm{stat}}^i + Z^\top H Y \Delta y_y, \\
    \tilde{r}_{\mathrm{ineq, redu}}^i &=  r_{\mathrm{ineq}}^i + G Y \Delta y_y - \Lambda_i^{-1}r_{\mathrm{comp}}^i.
\end{align*}
This formulation eliminates the equality constraints, recasting the optimal control structure into a general QP form, 
and yields a saddle-point system coupling the reduced primal variable $\Delta y_z$ (with regularized reduced Hessian $W_i$) and the dual variable $\Delta\lambda$.
Eliminating $\Delta\lambda$ by a further Schur complement recovers the complemented reduced Hessian of \eqref{eq:schur_comp_KKT_reduced} as:
\begin{equation*}
    \tilde{W}_i = W_i + Z^\top G^\top T_i^{-1}\Lambda_i G Z = \tilde{M}_i + E_i.
\end{equation*}

We will make the following assumptions for the proof:
\begin{assumption}
    \label{as:minimizer}
    $(y^*, \lambda^*)$ is a local (or global) minimizer for \eqref{eq:qp_problem}.
\end{assumption}

\begin{assumption}
    \label{as:isolated}
    Each local minimizer $(y^*,\lambda^*)$ for \eqref{eq:qp_problem} is isolated.
\end{assumption}

\begin{assumption}
    \label{as:licq}
    Linear independence constraint qualification (LICQ) holds at the minimizer $(y^*, \lambda^*)$.
\end{assumption}

\begin{assumption}
    \label{as:strict_comp}
    Strict complementarity holds at the minimizer $(y^*, \lambda^*)$.
\end{assumption}

Assumptions~\ref{as:minimizer}--\ref{as:licq} imply the second order sufficient condition (SOSC) for local optimality,
so we have the following Lemma~\ref{lem:SOSC}:
\begin{lemma}
    \label{lem:SOSC}
    Under Assumptions~\ref{as:minimizer}--\ref{as:licq}, for all $v \in C(y^*,\lambda^*) \setminus \{0\}$, it holds that $v^\top H v > 0$.
\end{lemma}

\begin{proof}
    With Assumption~\ref{as:minimizer}, the second order necessary condition (SONC) for local optimality holds, which states that for all $v \in C(y^*,\lambda^*)$, we have $v^\top H v \ge 0$.

    We proceed by contradiction. Suppose there exists $v \in C(y^*,\lambda^*) \setminus \{0\}$ such that $v^\top H v = 0$. 
    Consider the point $y(\alpha) = y^* + \alpha v$ for $\alpha > 0$. Because the constraints are linear and $v$ is in the critical cone, $y(\alpha)$ remains feasible for sufficiently small $\alpha$. 

    Since the objective function of \eqref{eq:qp_problem} is quadratic, expanding $f(y(\alpha))$ yields:
    \begin{equation*}
        f(y(\alpha)) = f(y^*) + \alpha \nabla f(y^*)^\top v + \frac{1}{2} \alpha^2 v^\top H v
    \end{equation*}

    As $\nabla f(y^*)^\top v = 0$, given the assumption that $v^\top H v = 0$, the equation reduces exactly to $f(y(\alpha)) = f(y^*)$. 
    This implies that $y(\alpha)$ is also a local minimizer, which directly contradicts Assumption~\ref{as:isolated}. Therefore, $v^\top H v > 0$ must hold.
\end{proof}

\begin{proposition}
    \label{prop:vanish}
    Under Assumptions~\ref{as:minimizer}--\ref{as:strict_comp}, given the sequence $\{\tilde{W}_i\}$ constructed by the Riccati-based regularization, as $\{(y_i, \lambda_i)\} \to (y^*, \lambda^*)$,
    $\tilde{W}_i = \tilde{M}_i = Z^\top \tilde{H}_i Z \succ 0$ for all $i$ sufficiently large.
\end{proposition}

\begin{proof}
    Under Assumption~\ref{as:strict_comp}, the critical cone $C(y^*,\lambda^*)$ reduces to the tangent space of the active constraints. 
    Therefore, we have $C(y^*,\lambda^*) = \{v \mid F v = 0,\ G_{\mathcal{A}} v = 0\}$.

    For any $v$ with $\|v\| = 1$, write $d = Zv \in \{d \mid Fd = 0\} \setminus \{0\}$
    and let $G_j$ denote the $j$-th row of $G$.
    The corresponding quadratic form of the reduced Hessian reads:
    \begin{align*}
        v^\top \tilde{M}_i v &= d^\top (H + G^\top T_i^{-1} \Lambda_i  G) d \\
                             &= d^\top H d + \sum_{j} \frac{\lambda_j^i}{t_j^i} (G_j d)^2 \\
                             &= d^\top H d + \sum_{j \in \mathcal{A}} \frac{\lambda_j^i}{t_j^i} (G_j d)^2 + \sum_{j \notin \mathcal{A}} \frac{\lambda_j^i}{t_j^i} (G_j d)^2
    \end{align*}

    We now proceed by contradiction. Assume there exists a sequence $v_i$ with $\|v_i\| = 1$ and, passing to a subsequence if necessary, $v_i \to \bar{v}$
    such that $v_i^\top \tilde{M}_i v_i \leq 0$ for all $i$.
    Writing $d_i = Z v_i$ and $\bar{d} = Z \bar{v}$, so that $d_i \to \bar{d}$, we can distinguish two cases:

    For $G_{\mathcal{A}}\bar{d}=0$, the limit direction satisfies $\bar{d} \in C(y^*,\lambda^*) \setminus \{0\}$,
    so $\bar{d}^\top H \bar{d} > 0$ follows by Lemma~\ref{lem:SOSC}.
    Since the remaining terms are non-negative, $v_i^\top \tilde{M}_i v_i \ge d_i^\top H d_i$ 
    with $d_i^\top H d_i \to \bar{d}^\top H \bar{d} > 0$.
    Hence $v_i^\top \tilde{M}_i v_i > 0$ for all sufficiently large $i$, contradicting the assumption.

    For $G_{\mathcal{A}}\bar{d} \neq 0$, Assumption~\ref{as:strict_comp} gives $\lambda_j^i \to \lambda_j^* > 0$ and $t_j^i \to 0$ for $j \in \mathcal{A}$,
    so the term $\sum_{j \in \mathcal{A}} \frac{\lambda_j^i}{t_j^i} (G_j d_i)^2 \rightarrow \infty$, 
    which dominates the other two terms, and thus $v_i^\top \tilde{M}_i v_i \rightarrow \infty$, contradicting the assumption.

    Therefore, we can conclude that $\tilde{W}_i = \tilde{M}_i = Z^\top \tilde{H}_i Z \succ 0$ for all $i$ sufficiently large, i.e., $E_i=0$.
\end{proof}

Proposition~\ref{prop:vanish} shows the regularization is asymptotically inactive: once the iterates are close to $(y^*,\lambda^*)$, $E_i=0$ and the regularized recursion coincides
with the unregularized one. The solver thus recovers the exact Newton step and inherits the local convergence rate of the underlying interior-point method, without the regularization affecting its asymptotic behavior.




% \section{Asymptotic Inactivity of the Regularization}
% \label{sec:asymptotic_inactivity}
% In this section, 
% we will prove that the reduced Hessian regularization via Riccati recursion would be asymptotically inactive,
% which means that the regularization will be inactive when the iterates are sufficiently close to the optimal solution $(y^*, \lambda^*)$,
% following a similar proof strategy to \cite{absilNewtonKKTInteriorpointMethods2007}.

% We will start with the regularized KKT formulation, where the null space method is applied to \eqref{eq:IPM_KKT_system} before the Schur complement.
% The resulting formulation is given as:
% \begin{equation}
%     \begin{bmatrix}
%     W & Z^\top G^\top\\
%     GZ & -\Lambda_i^{-1}T_i\\
%     \end{bmatrix}
%     \begin{bmatrix}
%     \Delta y_z \\
%     \Delta \lambda \\
%     \end{bmatrix} = -
%     \begin{bmatrix}
%     r_{\mathrm{stat, redu}}^i \\
%     \tilde{r}_{\mathrm{ineq, redu}}^i \\
%     \end{bmatrix},
%     \label{eq:reduced_KKT_system_schur}
% \end{equation}
% where
% \begin{align*}
%     W &= M + E_i, \\
%     r_{\mathrm{stat, redu}}^i &= Z^\top r_{\mathrm{stat}}^i - Z^\top H Y \Delta y_y, \\
%     \tilde{r}_{\mathrm{ineq, redu}}^i &=  r_{\mathrm{ineq}}^i - G Y \Delta y_y - \Lambda_i^{-1}r_{\mathrm{comp}}^i.
% \end{align*}
% This formulation eliminates the equality constraints, recasting the optimal control structure into a general QP form, 
% and yields a saddle-point system coupling the reduced primal variable $\Delta y_z$ (with regularized reduced Hessian $W$) and the dual variable $\Delta\lambda$.

% We will make the following assumptions for the proof:
% \begin{assumption}
%     \label{as:minimizer}
%     $(y^*, \lambda^*)$ is a local (or global) minimizer for \eqref{eq:qp_problem}.
% \end{assumption}

% \begin{assumption}
%     \label{as:isolated}
%     Each local minimizer $(y^*,\lambda^*)$ for \eqref{eq:qp_problem} is isolated.
% \end{assumption}

% \begin{assumption}
%     \label{as:licq}
%     Linear independence constraint qualification (LICQ) holds at the minimizer $(y^*, \lambda^*)$.
% \end{assumption}

% \begin{assumption}
%     \label{as:strict_comp}
%     Strict complementarity holds at the minimizer $(y^*, \lambda^*)$.
% \end{assumption}

% Assumptions~\ref{as:minimizer}--\ref{as:licq} imply the second order sufficient condition (SOSC) for local optimality,
% so we have the following Lemma~\ref{lem:SOSC}:
% \begin{lemma}
%     \label{lem:SOSC}
%     Under Assumptions~\ref{as:minimizer}--\ref{as:licq}, for all $v \in C(y^*,\lambda^*) \setminus \{0\}$, it holds that $v^\top H v > 0$.
% \end{lemma}

% \begin{proof}
%     With Assumption~\ref{as:minimizer}, the second order necessary condition (SONC) for local optimality holds, which states that for all $v \in C(y^*,\lambda^*)$, we have $v^\top H v \ge 0$.

%     We proceed by contradiction. Suppose there exists $v \in C(y^*,\lambda^*) \setminus \{0\}$ such that $v^\top H v = 0$. 
%     Consider the point $y(\alpha) = y^* + \alpha v$ for $\alpha > 0$. Because the constraints are linear and $v$ is in the critical cone, $y(\alpha)$ remains feasible for sufficiently small $\alpha$. 

%     Since the objective function of \eqref{eq:qp_problem} is quadratic, expanding $f(y(\alpha))$ yields:
%     \begin{equation*}
%         f(y(\alpha)) = f(y^*) + \alpha \nabla f(y^*)^\top v + \frac{1}{2} \alpha^2 v^\top H v
%     \end{equation*}

%     As $\nabla f(y^*)^\top v = 0$, given the assumption that $v^\top H v = 0$, the equation reduces exactly to $f(y(\alpha)) = f(y^*)$. 
%     This implies that $y(\alpha)$ is also a local minimizer, which directly contradicts Assumption~\ref{as:isolated}. Therefore, $v^\top H v > 0$ must hold.
% \end{proof}

% \begin{proposition}
%     \label{prop:vanish}
%     Under Assumptions~\ref{as:minimizer}--\ref{as:strict_comp}, given the sequence $\{W_i\}$ constructed by the algorithm, as $\{(y_i, \lambda_i)\} \to (y^*, \lambda^*)$,
%     $W_i = \tilde{M}_i = Z^\top \tilde{H}_i Z \succ 0$ for all $i$ sufficiently large.
% \end{proposition}

% \begin{proof}
%     Under Assumption~\ref{as:strict_comp}, the critical cone $C(y^*,\lambda^*)$ reduces to the tangent space of the active constraints. 
%     Therefore, we have $C(y^*,\lambda^*) = \{v \mid F v = 0 ,G_{\mathcal{A}} v = 0\}$.

%     For any $ v \in \{v \mid \|v\| = 1 \}$, we have the corresponding quadratic form of the reduced Hessian as:
%     \begin{align*}
%         v^\top \tilde{M}_i v &= d^\top (H + G^\top T_i^{-1} \Lambda_i  G) d \\
%                              &= d^\top H d + \sum_{j} \frac{\lambda_j}{s_j} (G_j d)^2 \\
%                              &= d^\top H d + \sum_{j \in \mathcal{A}} \frac{\lambda_j}{s_j} (G_j d)^2 + \sum_{j \notin \mathcal{A}} \frac{\lambda_j}{s_j} (G_j d)^2
%     \end{align*}
%     where $d = Zv \in \{d \mid Ad = 0\} \setminus \{0\}$. 

%     We now proceed by contradiction. Assume $\exists v_i \in \{v \mid \|v\| = 1 \}, v_i \rightarrow \bar{v} : v_i^\top \tilde{M}_i v_i \leq 0$.
%     We can distinguish two cases:

%     For $G_{\mathcal{A}}\bar{d}=0$, this indicates $d_i \in C(y^*,\lambda^*) \setminus \{0\}$. $d_i^\top H d_i > 0$ follows by Lemma~\ref{lem:SOSC}.
%     The rest of the terms are non-negative, so we have $\bar{v}^\top \tilde{M}_i \bar{v}> 0$, contradicting the assumption.

%     For $G_{\mathcal{A}}\bar{d} \neq 0$, the term $\sum_{j \in \mathcal{A}} \frac{\lambda_j}{s_j} (G_j d)^2 \rightarrow \infty$, 
%     which dominates the other two terms, and thus $v_i^\top \tilde{M}_i v_i \rightarrow \infty > 0$, contradicting the assumption.

%     Therefore, we can conclude that $W_i = \tilde{M}_i = Z^\top \tilde{H}_i Z \succ 0$ for all $i$ sufficiently large, i.e., $E_i=0$.
% \end{proof}

% Proposition~\ref{prop:vanish} shows the regularization is asymptotically inactive: once the iterates are close to $(y^*,\lambda^*)$, $E_i=0$ and the regularized recursion coincides
% with the unregularized one. The solver thus recovers the exact Newton step and inherits the local convergence rate of the underlying interior-point method, without the regularization affecting its asymptotic behavior.


% \subsection{Local Quadratic Convergence Rate for Quadratic Programs}
% \label{subsec:local_quadratic_convergence}
% In this section, 
% we will attempt to prove the local quadratic convergence of the regularized interior point method for solving indefinite QPs.
% Following similar proof structure as in \cite{absilNewtonKKTInteriorpointMethods2007}.

% We will start with the KKT formulation where the null space method is applied before the Schur complement to the \eqref{eq:IPM_KKT_system},
% the resulting formulation is given as:
% \begin{equation}
%     \begin{bmatrix}
%     M & Z^\top G^\top\\
%     GZ & -\Lambda_i^{-1}T_i\\
%     \end{bmatrix}
%     \begin{bmatrix}
%     \Delta y_z \\
%     \Delta \lambda \\
%     \end{bmatrix} = -
%     \begin{bmatrix}
%     r_{\mathrm{stat, redu}}^i \\
%     \tilde{r}_{\mathrm{ineq, redu}}^i \\
%     \end{bmatrix},
%     \label{eq:reduced_KKT_system_schur}
% \end{equation}
% where
% \begin{align*}
%     r_{\mathrm{stat, redu}}^i &= Z^\top r_{\mathrm{stat}}^i - Z^\top H Y \Delta y_y \\
%     \tilde{r}_{\mathrm{ineq, redu}}^i &=  r_{\mathrm{ineq}}^i - G Y \Delta y_y - \Lambda_i^{-1}r_{\mathrm{comp}}^i,
% \end{align*}
% this formulation eliminates the equality constraints, recasting the optimal control structure into a general QP form, 
% and yields a saddle-point system coupling the reduced primal variable $\Delta y_z$ (with reduced Hessian $M$) and the dual variable $\Delta\lambda$.

% We will make the following assumptions for the local quadratic convergence proof:
% \begin{assumption}
%     $(y^*, \lambda^*)$ is a local (or global) minimizer for \ref{eq:qp_problem}.
% \end{assumption}

% \begin{assumption}
%     Each local minimizer $(y^*,\lambda^*)$ for \ref{eq:qp_problem} is isolated.
% \end{assumption}

% \begin{assumption}
%     Linear independence constraint qualification (LICQ) holds at the minimizer $(y^*, \lambda^*)$.
% \end{assumption}

% \begin{assumption}
%     Strict complementarity holds at the minimizer $(y^*, \lambda^*)$.
% \end{assumption}

% The Assumptions 1-3 imply the second order sufficient condition (SOSC) for local optimality,
% so we have the following lemma 1:
% \begin{lemma}
%     Under Assumptions 1-3, for all $v \in C(y^*,\lambda^*) \setminus \{0\}$, it holds that $v^\top H v > 0$.
% \end{lemma}

% \begin{proof}
%     With Assumption 1, The second order necessary condition (SONC) for local optimality holds, which states that for all $v \in C(y^*,\lambda^*)$, we have $v^\top H v \ge 0$.

%     We proceed by contradiction. Suppose there exists $v \in C(y^*,\lambda^*) \setminus \{0\}$ such that $v^\top H v = 0$. 
%     Consider the point $y(\alpha) = y^* + \alpha v$ for $\alpha > 0$. Because the constraints are linear and $v$ is in the critical cone, $y(\alpha)$ remains feasible for sufficiently small $\alpha$. 

%     Since the objective function of \ref{eq:qp_problem} is quadratic. Expanding $f(y(\alpha))$ yields:
%     $$f(y(\alpha)) = f(y^*) + \alpha \nabla f(y^*)^\top v + \frac{1}{2} \alpha^2 v^\top H v$$

%     As $\nabla f(y^*)^\top v = 0$, given the assumption that $v^\top H v = 0$, the equation reduces exactly to $f(y(\alpha)) = f(y^*)$. 
%     This implies that $y(\alpha)$ is also a local minimizer, which directly contradicts Assumption 2. Therefore, $v^\top H v > 0$ must hold.
%     \label{proof:SOSC}
% \end{proof}

% \begin{lemma}
%     Under Assumption 1-4, given the sequence $\{W_i\}$ constructed by the algorithm, as \{$(y_i, \lambda_i)$\} $\rightarrow$ $(y^*, \lambda^*)$,
%     $W_i = \tilde{M}_i = Z^\top \tilde{H}_i Z \succ 0$ for all $i$ sufficiently large.
% \end{lemma}

% \begin{proof}
%     Under Assumption 4, the critical cone $C(y^*,\lambda^*)$ reduces to the tangent space of the active constraints. 
%     Therefore, we have $C(y^*,\lambda^*) = \{v \mid F v = 0 ,G_{\mathcal{A}} v = 0\}$.

%     We then proceed with the reduced Hessian. For any $ v \in \{v \mid \|v\| = 1 \}$, we have:
%     \begin{align*}
%         v^\top \tilde{M}_i v &= d^\top (H + G^\top T_i^{-1} \Lambda_i  G) d \\
%                              &= d^\top H d + \sum_{j} \frac{\lambda_j}{s_j} (G_j d)^2 \\
%                              &= d^\top H d + \sum_{j \in \mathcal{A}} \frac{\lambda_j}{s_j} (G_j d)^2 + \sum_{j \notin \mathcal{A}} \frac{\lambda_j}{s_j} (G_j d)^2
%     \end{align*}
%     where $d = Zv \in \{d \mid Ad = 0\} \setminus \{0\}$. 

%     We now proceed by contradiction, Assume $\exists v_i \in \{v \mid \|v\| = 1 \}, v_i \rightarrow \bar{v} : v_i^\top \tilde{M}_i v_i \leq 0$.
%     We can distinguish two cases:

%     For $G_{\mathcal{A}}\bar{d}=0$, this indicates $d_i \in C(y^*,\lambda^*) \setminus \{0\}$. $d_i^\top H d_i > 0$ follows by Lemma 1.
%     Rest of the terms are non-negative, so we have $\bar{v}^\top \tilde{M}_i \bar{v}> 0$, contradict with assumption.

%     For $G_{\mathcal{A}}\bar{d} \neq 0$, the term $\sum_{j \in \mathcal{A}} \frac{\lambda_j}{s_j} (G_j d)^2 \rightarrow \infty$, 
%     which dominates the other two terms, and thus $v_i^\top \tilde{M}_i v_i \rightarrow \infty > 0$, contradict with assumption.

%     Therefore, we can conclude that $W_i = \tilde{M}_i = Z^\top \tilde{H}_i Z \succ 0$ for all $i$ sufficiently large, i.e., $E_i=0$.
% \end{proof}


% We take part of the following proposition from \cite{titsSimpleQuadraticallyConvergent1994} as a key step for proof.

% \begin{proposition}
% Let $F : \mathcal{R}^n \to \mathcal{R}^n$ be twice continuously differentiable and let $w^* \in \mathcal{R}^n$.
% Let $\rho > 0$ be such that $F(w^*) = 0$ and $\frac{\partial F}{\partial w}(w)$ is invertible whenever $w \in \overline{B}(w^*, \rho)$.
% Let $\delta^N : \overline{B}(w^*, \rho) \to \mathcal{R}^n$ be the Newton increment $\delta^N(w) := -\left(\frac{\partial F}{\partial w}(w)\right)^{-1} F(w)$.
% Then given any $c_1 > 0$ there exist $c_2 > 0$ such that the following statement holds:

% For every $w \in \overline{B}(w^*, \rho)$ and every $w^+ \in \mathcal{R}^n$ satisfying
% \begin{equation}
%     \|w^+ - (w + \delta^N(w))\| \leq c_1 \|\delta^N(w)\|^2,
%     \label{ineq:condition}
% \end{equation}
% it holds that
% \begin{equation}
%     \|w^+ - w^*\| \leq c_2 \|w - w^*\|^2.
% \end{equation}
% \end{proposition}

% With the Proposition 3, we will define $w = (y, \lambda)$ as the primal-dual variable and $w^* = (y^*, \lambda^*)$ as the solution,
% and establish the local quadratic convergence by showing that the conditions in Proposition 3 hold for the iterates generated by the algorithm.

% \begin{lemma}(Nocedal).
%     Under Assumption 4, when \{$(y_i, \lambda_i)$\} $\rightarrow$ $(y^*, \lambda^*)$, $\sigma_i$ $\mu_i$ = O($||\delta^N||^4$).
% \end{lemma}

% \begin{proof}
%     By Assumption 4, we have:
%     \begin{equation*}
%         \mu = \frac{1}{m} \sum_i t_i \lambda_i = O(\|w - w^*\|) = O(\|\delta^N\|).
%     \end{equation*}
%     Since the affine-scaling step is the Newton step for $F(w) = 0$, we obtain:
%     \begin{align*}
%         t_i + \Delta t_i^{\mathrm{aff}} &= t_i^* + O(\|w - w^*\|^2), \\
%         \lambda_i + \Delta \lambda_i^{\mathrm{aff}} &= \lambda_i^* + O(\|w - w^*\|^2).
%     \end{align*}
%     It follows that $(t_i + \Delta t_i^{\mathrm{aff}})(\lambda_i + \Delta \lambda_i^{\mathrm{aff}}) = O(\|w - w^*\|^2)$. 
%     Thus, $\mu^{\mathrm{aff}} = O(\|w - w^*\|^2) = O(\mu^2)$. 

%     According to the update rule, we obtain $\sigma = (\mu^{\mathrm{aff}} / \mu)^3 = O(\mu^3) \to 0$. 
%     Therefore, combining the orders derived above yields:
%     \begin{equation*}
%         \sigma\mu = O(\mu^4) = O(\|\delta^N\|^4).
%     \end{equation*}
% \end{proof}

% \begin{lemma}
%     Under Assumption 4, when \{$(y_i, \lambda_i)$\} $\rightarrow$ $(y^*, \lambda^*)$, the affine step size of primal step and dual step satisfy
%     \begin{equation*}
%         |1-\alpha_p| = O(\|\delta^N\|), \qquad |1-\alpha_d| = O(\|\delta^N\|).
%     \end{equation*}
% \end{lemma}

% \begin{proof}
% For the affine step, the linearized complementary condition is given by:
% \begin{equation*}
% t_i \Delta \lambda_i^{\mathrm{aff}} + \lambda_i \Delta t_i^{\mathrm{aff}} = -t_i \lambda_i
% \end{equation*}

% Dividing both sides by $t_i \lambda_i$ yields the fundamental step identity:
% \begin{equation*}
% \frac{\Delta \lambda_i^{\mathrm{aff}}}{\lambda_i} + \frac{\Delta t_i^{\mathrm{aff}}}{t_i} = -1,
% \end{equation*}
% where the magnitude bounds for the affine step are given from Lemma 4 as:
% \begin{equation*}
%     \|\Delta t^{\mathrm{aff}}\| = O(\|\delta^N\|), \qquad \|\Delta \lambda^{\mathrm{aff}}\| = O(\|\delta^N\|).
% \end{equation*}

% \emph{Primal step.}
% As the FTB for primal variables is defined as:
% \begin{equation*}
%     \alpha_p^{\max} = \min\left(1, \min_{\Delta t_j < 0} \frac{t_j}{-\Delta t_j}\right),
% \end{equation*} 
% where based on the identity, the ratio $\frac{t_j}{-\Delta t_j}$ can be expressed as:
% \begin{equation*}
%     \frac{t_j}{-\Delta t_j} = \frac{1}{1 + \frac{\Delta \lambda_i}{\lambda_i}}.
% \end{equation*}

% For inactive constraints $t_j\rightarrow t^\star > 0$, thus the ratio $\frac{t_j}{-\Delta t_j}  \rightarrow \infty$, and such indices do not attain the minimum.

% Since $\alpha_p^{\max} \leq 1$, for active constraints, $\lambda_i \rightarrow \lambda_i^* > 0$ 
% and $\Delta \lambda_i = O(\|\delta^N\|)$, we have:
% \begin{equation*}
%     1 - \alpha_p^{\max} = 1 - \frac{t_j}{-\Delta t_j} = \frac{\Delta \lambda}{\lambda + \Delta \lambda} = O(\|\delta^N\|).
% \end{equation*}

% Considering the safeguard value:
% \begin{align*}
%     \Delta(\alpha_p^{\max}) &= (1-\alpha_p^{\max}) \epsilon_\alpha + (1 - \mu_c) \alpha_p^{\max},
% \end{align*}
% where $\mu_c = \min(\epsilon_\mu, \mu)$.

% We have:
% \begin{align*}
%     \alpha_p &= \tau(\alpha_p^{\max}) \alpha_p^{\max} \\
%              &= \alpha_p^{\max}[1-(1-\alpha_p^{\max})(1-\epsilon_\alpha) - \mu_c\alpha_p^{\max}] \\
%              &= \alpha_p^{\max} - (1-\alpha_p^{\max})(1-\epsilon_\alpha)\alpha_p^{\max} - \mu_c(\alpha_p^{\max})^2
% \end{align*}

% In local convergence, $\mu_c = \mu$ holds, by lemma 4, the $1 - \alpha_p$ is :
% \begin{align*}
%     1 - \alpha_p &= O(\|\delta^N\|) + O(\|\delta^N\|) + O(\|\delta^N\|^4) \\
%                  &= O(\|\delta^N\|).
% \end{align*}

% \emph{Dual step.}
% Symmetrically, the FTB for dual variables is defined as:
% \begin{equation*}
%     \alpha_d^{\max} = \min\left(1, \min_{\Delta \lambda_j < 0} \frac{\lambda_j}{-\Delta \lambda_j}\right),
% \end{equation*}
% For active constraints, corresponding indices do not attain the minimum, due to $\frac{\lambda}{-\Delta \lambda} \rightarrow \infty$.
% For inactive constraints, $t_j \rightarrow t_j^* > 0$ and $\Delta t_j = O(\|\delta^N\|)$, 
% thus the ratio $\frac{\lambda_j}{-\Delta \lambda_j} = 1 - O(\|\delta^N \|)$.

% We then have:
% \begin{equation*}
%     1 - \alpha_d^{\max} =  1 - \frac{\lambda_j}{-\Delta \lambda_j} = O(\|\delta^N\|).
% \end{equation*}

% Following the same logic as primal step, we can conclude that $1 - \alpha_d = O(\|\delta^N\|)$.

% Combining the two estimates gives 
% \begin{equation*}
%     |1-\alpha_p| = O(\|\delta^N\|), \qquad |1-\alpha_d| = O(\|\delta^N\|).
% \end{equation*}
% \end{proof}

% \begin{theorem}
%     Under Assumption 1-4, Let $\{y_i\}$, $\{\lambda_i\}$ as sequences constructed by algorithm, 
%     $\{(y_i, \lambda_i)\}$ converges to the solution $(y^*, \lambda^*)$ with local quadratic convergence rate.
% \end{theorem}

% \begin{proof}
%     We establish that the conditions in Proposition 3 hold with \ref{eq:reduced_KKT_system_schur} as:

%     \begin{equation}
%         \underbrace{
%             \begin{bmatrix}
%             M + E_i & Z^\top G^\top\\
%             GZ & -\Lambda_i^{-1}T_i\\
%             \end{bmatrix}
%             \begin{bmatrix}
%             \Delta y_z \\
%             \Delta \lambda \\
%             \end{bmatrix}
%         }_{\mathcal{M}(w)} = -
%         \begin{bmatrix}
%         r_{\mathrm{stat, redu}}^i \\
%         \tilde{r}_{\mathrm{ineq, redu}}^i \\
%         \end{bmatrix},
%     \end{equation}

%     Under Assumptions 1-4, $\frac{\partial \mathcal{M}}{\partial w}$ is invertible at $w^*$.
%     By Lemma 2, $E_i=0$ for all $i$ sufficiently large,
%     so the left-hand side coincides with the exact Newton Jacobian.
%     By Lemma 4,
%     $\sigma_i\mu_i = O(\|\delta^N\|^4)$. Consequently,$||\Delta w_i - \delta^N(w_i)|| = O(\|\delta^N(w^i)\|^4)$ for $w^i \to w^*$.

%     With $\Delta w_i = [\Delta y_i, \Delta \lambda_i]^\top$, we start with the update of primal variables $y$,
%     \begin{align}
%         \|y_{i+1} - (y_i+ \delta^N_y)\| &=  \|y_{i}+ \alpha_p \Delta y_i - (y_i+ \delta^N_y)\| \\
%                                       &\le |1 - \alpha_p| ||\Delta y_i|| + ||\Delta y_i - \delta^N_y||\\
%                                       &\le C_1 ||\delta^N_y|| ||\Delta y_i|| + C_2 ||\delta^N_y||^4 \\
%                                       &\le C' ||\delta^N_y||^2,
%     \end{align}
%     As a result, the condition \ref{ineq:condition} in Proposition 3 holds for the primal variables.
%     For the dual variables $\lambda$, the same logic applies with $\alpha_d$ substituting $\alpha_p$.
%     \footnote{In the implementation, 
%     the update for the dual variables is taken as $\max(\lambda+\alpha_d\Delta\lambda, \epsilon)$, 
%     where $\epsilon$ is on the order of machine precision.
%     This truncation is only activated when the accuracy reaches $\mathcal{O}(\sqrt{\epsilon})$ and thus does not affect the asymptotic convergence rate analysis above.}

%     Thus, we can conclude that the sequence $\{(y_i, \lambda_i)\}$ converges to $(y^*, \lambda^*)$ with local quadratic convergence rate.
% \end{proof}

\end{document}
